\section{Comparison with the native core}
\label{mir:comparison}


The safe equality operations, NC, and EI are now available with their
full contextual force. The zero-source restoration argument in
Theorem~\ref{compat:zero-restored} therefore applies: it derives
$0=t\vdash t=0$, full numerical equality algebra, and unrestricted
canonical formula transport. Its only nontrivial induction uses the
equation $n=I(n,t;u.Su)$ with $t$ fixed; Theorem~\ref{mir:equation-induction}
supplies precisely that rule without a forbidden proof substitution.
Case DI for every canonical matrix was already obtained in
Lemma~\ref{mir:canonical-cases}, so the binder step in the restoration
of general formula transport is legitimate.

Let $\mathcal C$ be the inherited Choice-DI calculus: the same terms,
all formulas generated freely from equality and $D$, the same
computation, Use and DE rules, unrestricted transport, case DI at
every matrix, and progressive DI at literal equations. Define
\[
 T(s=t):=(s=t),\qquad
 T(D_{x,y}B):=D_{x,y}P(TB),
 \tag{Translation}\label{mir:translation}
\]
renaming all padding variables freshly. Translate a context clause by
$T\langle A;y.B\rangle:=\langle TA;y.TB\rangle$. Let
$J:\mathcal M\to\mathcal C$ be literal inclusion of the canonical
syntax. Neither translation changes a numerical term.

\begin{theorem}[Translation of derivations]\label{mir:translation-theorem}
The following effective transformations preserve the full consequence
relations:
\[
 \Gamma\vdash_{\mathcal C}A\Longrightarrow
 T\Gamma\vdash_{\mathcal M}TA,
 \qquad
 \Gamma\vdash_{\mathcal M}B\Longrightarrow
 J\Gamma\vdash_{\mathcal C}JB.
\]
\end{theorem}
\begin{proof}
Inclusion preserves every primitive rule: source-restricted transport
is an unrestricted instance, and root-directed DI chooses an available
case or progress rule. For $T$, induct on a $\mathcal C$ derivation.
Computation is unchanged. Use commutes with formula substitution.
Transport uses the translated template and restored unrestricted
transport.

A case-DI step is translated by introducing the compound matrix
$P(TB)$. Unpad its local base, apply the translated premise, and pad
the successor result. For a progressive step at an equation $E$,
the target is $D_{x,y}L(E)$. Introduce it by compound case DI.
Rename the supplied step eigenparameter freshly, away from the new
outer successor target, before applying EI.
Its base $L(E(t,0))$ supplies $E(t,0)$. Recovered EI applied to the
translated step then gives $E(t,Sz)$, which padding converts into
the required successor target. Before this EI application, rename
the step's eigenparameter to a fresh $w$, distinct from the outer
case parameter $z$ occurring in the target $Sz$. This is a safe
renaming, not substitution of a numerical term into a derivation.

For DE, the translated major opens the clause
$\langle P(TB)(b,0);y.P(TB)(b,y)\rangle$. The translated minor
expects $\langle TB(b,0);y.TB(b,y)\rangle$. Replace each use of
that expected clause by padding its antecedent, using the available
clause at the same target, and unpadding its consequent. This local
replacement preserves the fresh-witness condition. All other cases
are immediate.
\end{proof}

\begin{lemma}[Canonical atomic composite]\label{mir:atomic-composite}
For every admissible mirror matrix $B(x,y)$,
\[
              D_{x,y}B\quad\dashv\vdash_{\mathcal M}\quad
              D_{x,y}L(B).
\]
\end{lemma}
\begin{proof}
Forwards, open the source clause at a fresh witness $b$ and introduce
the compound target matrix by case DI. Unpad its base, use the
source clause at $Sz$, and pad its successor instance.

Backwards, open the clause of $D_{x,y}L(B)$ at $b$. The mirror
zero-instance theorem proves $B(b,0)$, hence $L(B(b,0))$. The
opened clause yields $L(B(b,Sz))$, and unpadding gives $B(b,Sz)$
independently of the local hypothesis $B(b,z)$. Atomic progressive
DI therefore introduces $D_{x,y}B$ at witness $b$. Close DE.
\end{proof}

\begin{theorem}[Both composites and the whole-domain interpretation]
\label{mir:bi-interpretation}
For every $\mathcal C$-formula $A$ and canonical $\mathcal M$-formula $B$,
\[
 A\dashv\vdash_{\mathcal C}JT(A),\qquad
 B\dashv\vdash_{\mathcal M}TJ(B).
\]
The ordinary realization of $\mathcal M$ is bi-interpretable and
equiconsistent with iPRA through the inherited native-core bridge.
Both numerical domains are whole, numerical equality remains actual
equality, and both numerical composites have identity graph $x=y$.
\end{theorem}
\begin{proof}
Padding is interderivable with its argument in $\mathcal C$, and
case DI gives congruence under $D$. Structural induction therefore
gives $A\dashv\vdash JT(A)$.

In $\mathcal M$, uniformly interderivable compound matrices may be
exchanged under $D$: eliminate the source binder, introduce the
target by compound case DI, convert its base to the source base,
use the opened clause at a successor, and convert back. The reverse
direction is identical. Structural induction on a canonical formula
now proves the second composite. At a compound matrix use this
congruence; at an atomic mirror matrix use
Lemma~\ref{mir:atomic-composite}. The transformations are uniform in
retained contexts and therefore extend to clauses by replacing their
Use instances.

The ordinary reading leaves equations and terms unchanged and reads
\[
 D_{x,y}B\quad\text{as}\quad
 \exists x\bigl(B(x,0)\to\forall y\,B(x,y)\bigr),
\]
and $\langle A;y.B\rangle$ as $A\to\forall yB$. The numerical
interpretations in the two displayed native translations are literal
identities. Composing them, and their proved formula composites,
with the native-core/ordinary-scalar/iPRA interpretations gives the
claimed whole-domain bi-interpretation. The contradiction $0=1$
is unchanged, so equiconsistency follows. No standard-numeral
restriction, quotient equality, or uniform evaluator for function
descriptions is introduced.
\end{proof}

\paragraph{Scope of the reduction.}
The new primitive inventory uses one root-determined introduction and
omits all literal-zero-source transports. Its guard uses a single
iterator and repeats its initializer, on the opposite side of the
equation. The proofs above establish preservation of the full
consequence relation after the displayed formula translations. They
do not establish global minimality of an axiomatization or a shortest
possible implementation. A raw certificate checker that accepts
noncanonical goals must additionally state its administrative reading;
the theorem about canonical formulas alone is not a theorem that such
a program validates the canonical grammar.
